
\documentclass{article}
\usepackage{colortbl}
\usepackage{makecell}
\usepackage{multirow}
\usepackage{supertabular}

\begin{document}

\newcounter{utterance}

\centering \large Interaction Transcript for game `dond', experiment `coop\_en', episode 14 with reflexion{-}Qwen3.5{-}9B.
\vspace{24pt}

{ \footnotesize  \setcounter{utterance}{1}
\setlength{\tabcolsep}{0pt}
\begin{supertabular}{c@{$\;$}|p{.15\linewidth}@{}p{.15\linewidth}p{.15\linewidth}p{.15\linewidth}p{.15\linewidth}p{.15\linewidth}}
    \# & $\;$A & \multicolumn{4}{c}{Game Master} & $\;\:$B\\
    \hline

    \theutterance \stepcounter{utterance}  
    & & \multicolumn{4}{p{0.6\linewidth}}{
        \cellcolor[rgb]{0.9,0.9,0.9}{
            \makecell[{{p{\linewidth}}}]{
                \texttt{\tiny{[P1$\langle$GM]}}
                \texttt{You are playing a cooperative negotiation game in which you have to agree on how to divide a set of items among you and another player.} \\
\\ 
\texttt{Rules:} \\
\texttt{(a) You and the other player are given a set of items. Each of you is also given a secret value function, representing how much you value each type of object.} \\
\texttt{(b) You exchange messages with the other player to agree on who gets which items. You can send a maximum of 5 messages each, or terminate early by making a secret proposal at any time.} \\
\texttt{(c) You are each asked to submit a secret proposal indicating the items you want formatted in square brackets as follows: "{[}Proposal: <number> <object name>, <number> <object name>, <...>{]}"} \\
\texttt{(d) If your proposals are complementary, i.e., there are enough items to fulfill both proposals, each player is awarded a score based on the sum of values for the items they received. Otherwise, both of you get zero points.} \\
\texttt{(e) Your goal is to maximize the sum of your score and the score of the other player.} \\
\\ 
\texttt{Let us start.} \\
\\ 
\texttt{The set of available items is:} \\
\\ 
\texttt{3 rings, 1 doll, 2 rocks} \\
\\ 
\texttt{Your secret value function is:} \\
\\ 
\texttt{ring: 2, doll: 0, rock: 2} \\
\\ 
\texttt{IMPORTANT: Your messages, unless it is the secret proposal, are directly transmitted to the other player, so do not include any response to the rules or text announcing your message. To make a secret proposal, use the indicated format. Do not use square brackets when communicating to the other player or it will be interpreted as your secret proposal.} \\
            }
        }
    }
    & \\ \\

    \theutterance \stepcounter{utterance}  
    & \multicolumn{4}{p{0.6\linewidth}}{
        \cellcolor[rgb]{0.9,0.9,0.9}{
            \makecell[{{p{\linewidth}}}]{
                \texttt{\tiny{[P1$\rangle$GM]}}
                \texttt{I understand the rules and my valuation function. I am ready to negotiate.} \\
\\ 
\texttt{My valuations:} \\
\texttt{{-} Ring: 2} \\
\texttt{{-} Doll: 0} \\
\texttt{{-} Rock: 2} \\
\\ 
\texttt{Total value of all items: \$3 \textbackslash{}times 2 + 1 \textbackslash{}times 0 + 2 \textbackslash{}times 2 = 6 + 0 + 4 = 10\$.} \\
\\ 
\texttt{To maximize the sum of our scores (which is simply the total value of all items if we succeed), we need to ensure our proposals are complementary. Since I value the Doll at 0, I am indifferent to who gets it, but I prefer to get the Rings and Rocks.} \\
\\ 
\texttt{I will start by proposing a split that gives me the items I value most and leaves the rest for you, ensuring we can fulfill both requests.} \\
\\ 
\texttt{Proposal:} \\
\texttt{I want 2 rings and 2 rocks.} \\
\texttt{You get 1 ring and 1 doll.} \\
\\ 
\texttt{This uses:} \\
\texttt{{-} Rings: 2 (me) + 1 (you) = 3 (Available: 3) ✓} \\
\texttt{{-} Dolls: 0 (me) + 1 (you) = 1 (Available: 1) ✓} \\
\texttt{{-} Rocks: 2 (me) + 0 (you) = 2 (Available: 2) ✓} \\
\\ 
\texttt{The proposals are complementary.} \\
\\ 
\texttt{Message to you:} \\
\texttt{I propose taking 2 rings and 2 rocks. You can take the remaining 1 ring and the 1 doll. This way we use all items and both of us get items we value.} \\
            }
        }
    }
    & & \\ \\

    \theutterance \stepcounter{utterance}  
    & & & \multicolumn{2}{p{0.3\linewidth}}{
        \cellcolor[rgb]{0.9,0.9,0.9}{
            \makecell[{{p{\linewidth}}}]{
                \texttt{\tiny{[GM$|$GM]}}
                \texttt{continue} \\
            }
        }
    }
    & & \\ \\

    \theutterance \stepcounter{utterance}  
    & & & \multicolumn{2}{p{0.3\linewidth}}{
        \cellcolor[rgb]{0.9,0.9,0.9}{
            \makecell[{{p{\linewidth}}}]{
                \texttt{\tiny{[GM$|$GM]}}
                \texttt{I understand the rules and my valuation function. I am ready to negotiate.} \\
\\ 
\texttt{My valuations:} \\
\texttt{{-} Ring: 2} \\
\texttt{{-} Doll: 0} \\
\texttt{{-} Rock: 2} \\
\\ 
\texttt{Total value of all items: \$3 \textbackslash{}times 2 + 1 \textbackslash{}times 0 + 2 \textbackslash{}times 2 = 6 + 0 + 4 = 10\$.} \\
\\ 
\texttt{To maximize the sum of our scores (which is simply the total value of all items if we succeed), we need to ensure our proposals are complementary. Since I value the Doll at 0, I am indifferent to who gets it, but I prefer to get the Rings and Rocks.} \\
\\ 
\texttt{I will start by proposing a split that gives me the items I value most and leaves the rest for you, ensuring we can fulfill both requests.} \\
\\ 
\texttt{Proposal:} \\
\texttt{I want 2 rings and 2 rocks.} \\
\texttt{You get 1 ring and 1 doll.} \\
\\ 
\texttt{This uses:} \\
\texttt{{-} Rings: 2 (me) + 1 (you) = 3 (Available: 3) ✓} \\
\texttt{{-} Dolls: 0 (me) + 1 (you) = 1 (Available: 1) ✓} \\
\texttt{{-} Rocks: 2 (me) + 0 (you) = 2 (Available: 2) ✓} \\
\\ 
\texttt{The proposals are complementary.} \\
\\ 
\texttt{Message to you:} \\
\texttt{I propose taking 2 rings and 2 rocks. You can take the remaining 1 ring and the 1 doll. This way we use all items and both of us get items we value.} \\
            }
        }
    }
    & & \\ \\

    \theutterance \stepcounter{utterance}  
    & & \multicolumn{4}{p{0.6\linewidth}}{
        \cellcolor[rgb]{0.9,0.9,0.9}{
            \makecell[{{p{\linewidth}}}]{
                \texttt{\tiny{[GM$\rangle$P2]}}
                \texttt{You are playing a cooperative negotiation game in which you have to agree on how to divide a set of items among you and another player.} \\
\\ 
\texttt{Rules:} \\
\texttt{(a) You and the other player are given a set of items. Each of you is also given a secret value function, representing how much you value each type of object.} \\
\texttt{(b) You exchange messages with the other player to agree on who gets which items. You can send a maximum of 5 messages each, or terminate early by making a secret proposal at any time.} \\
\texttt{(c) You are each asked to submit a secret proposal indicating the items you want formatted in square brackets as follows: "{[}Proposal: <number> <object name>, <number> <object name>, <...>{]}"} \\
\texttt{(d) If your proposals are complementary, i.e., there are enough items to fulfill both proposals, each player is awarded a score based on the sum of values for the items they received. Otherwise, both of you get zero points.} \\
\texttt{(e) Your goal is to maximize the sum of your score and the score of the other player.} \\
\\ 
\texttt{Let us start.} \\
\\ 
\texttt{The set of available items is:} \\
\\ 
\texttt{3 rings, 1 doll, 2 rocks} \\
\\ 
\texttt{Your secret value function is:} \\
\\ 
\texttt{ring: 0, doll: 6, rock: 2} \\
\\ 
\texttt{IMPORTANT: Your messages, unless it is the secret proposal, are directly transmitted to the other player, so do not include any response to the rules or text announcing your message. To make a secret proposal, use the indicated format. Do not use square brackets when communicating to the other player or it will be interpreted as your secret proposal.} \\
\\ 
\\ 
\texttt{I understand the rules and my valuation function. I am ready to negotiate.} \\
\\ 
\texttt{My valuations:} \\
\texttt{{-} Ring: 2} \\
\texttt{{-} Doll: 0} \\
\texttt{{-} Rock: 2} \\
\\ 
\texttt{Total value of all items: \$3 \textbackslash{}times 2 + 1 \textbackslash{}times 0 + 2 \textbackslash{}times 2 = 6 + 0 + 4 = 10\$.} \\
\\ 
\texttt{To maximize the sum of our scores (which is simply the total value of all items if we succeed), we need to ensure our proposals are complementary. Since I value the Doll at 0, I am indifferent to who gets it, but I prefer to get the Rings and Rocks.} \\
\\ 
\texttt{I will start by proposing a split that gives me the items I value most and leaves the rest for you, ensuring we can fulfill both requests.} \\
\\ 
\texttt{Proposal:} \\
\texttt{I want 2 rings and 2 rocks.} \\
\texttt{You get 1 ring and 1 doll.} \\
\\ 
\texttt{This uses:} \\
\texttt{{-} Rings: 2 (me) + 1 (you) = 3 (Available: 3) ✓} \\
\texttt{{-} Dolls: 0 (me) + 1 (you) = 1 (Available: 1) ✓} \\
\texttt{{-} Rocks: 2 (me) + 0 (you) = 2 (Available: 2) ✓} \\
\\ 
\texttt{The proposals are complementary.} \\
\\ 
\texttt{Message to you:} \\
\texttt{I propose taking 2 rings and 2 rocks. You can take the remaining 1 ring and the 1 doll. This way we use all items and both of us get items we value.} \\
            }
        }
    }
    & \\ \\

    \theutterance \stepcounter{utterance}  
    & & & \multicolumn{4}{p{0.6\linewidth}}{
        \cellcolor[rgb]{0.9,0.9,0.9}{
            \makecell[{{p{\linewidth}}}]{
                \texttt{\tiny{[GM$\langle$P2]}}
                \texttt{I propose taking 1 ring and 1 rock. You can take the remaining 2 rings, 1 doll, and 1 rock. This way we use all items and both of us get items we value.} \\
            }
        }
    }
     \\ \\

    \theutterance \stepcounter{utterance}  
    & & & \multicolumn{2}{p{0.3\linewidth}}{
        \cellcolor[rgb]{0.9,0.9,0.9}{
            \makecell[{{p{\linewidth}}}]{
                \texttt{\tiny{[GM$|$GM]}}
                \texttt{continue} \\
            }
        }
    }
    & & \\ \\

    \theutterance \stepcounter{utterance}  
    & & & \multicolumn{2}{p{0.3\linewidth}}{
        \cellcolor[rgb]{0.9,0.9,0.9}{
            \makecell[{{p{\linewidth}}}]{
                \texttt{\tiny{[GM$|$GM]}}
                \texttt{I propose taking 1 ring and 1 rock. You can take the remaining 2 rings, 1 doll, and 1 rock. This way we use all items and both of us get items we value.} \\
            }
        }
    }
    & & \\ \\

    \theutterance \stepcounter{utterance}  
    & & \multicolumn{4}{p{0.6\linewidth}}{
        \cellcolor[rgb]{0.9,0.9,0.9}{
            \makecell[{{p{\linewidth}}}]{
                \texttt{\tiny{[P1$\langle$GM]}}
                \texttt{I propose taking 1 ring and 1 rock. You can take the remaining 2 rings, 1 doll, and 1 rock. This way we use all items and both of us get items we value.} \\
            }
        }
    }
    & \\ \\

    \theutterance \stepcounter{utterance}  
    & \multicolumn{4}{p{0.6\linewidth}}{
        \cellcolor[rgb]{0.9,0.9,0.9}{
            \makecell[{{p{\linewidth}}}]{
                \texttt{\tiny{[P1$\rangle$GM]}}
                \texttt{My proposal: {[}Proposal: 2 ring, 2 rock{]}} \\
\\ 
\texttt{Let's check if our proposals are complementary:} \\
\texttt{{-} **Rings**: I want 2, you want 1. Total needed: 3. Available: 3. (OK)} \\
\texttt{{-} **Dolls**: I want 0, you want 1. Total needed: 1. Available: 1. (OK)} \\
\texttt{{-} **Rocks**: I want 2, you want 1. Total needed: 3. Available: 2. (**Conflict**)} \\
\\ 
\texttt{We have a conflict with the rocks. I requested 2 rocks, but there are only 2 available in total, and you requested 1. This means we cannot fulfill both proposals as stated. If we proceed with these specific numbers, both of us get zero points.} \\
\\ 
\texttt{To maximize the sum of our scores (which is the total value of all items, 10), we need to adjust our requests so they fit within the available inventory.} \\
\\ 
\texttt{Since I value rocks at 2 and you also value them, and there are only 2 rocks total, we must split them.} \\
\texttt{{-} I will take 1 rock.} \\
\texttt{{-} You take 1 rock.} \\
\\ 
\texttt{Let's revise my proposal to be compatible with yours:} \\
\texttt{{-} **Rings**: I take 2, you take 1. (Total 3 {-} OK)} \\
\texttt{{-} **Dolls**: I take 0, you take 1. (Total 1 {-} OK)} \\
\texttt{{-} **Rocks**: I take 1, you take 1. (Total 2 {-} OK)} \\
\\ 
\texttt{This new set of requests uses all items perfectly.} \\
\texttt{{-} My score: 2 rings (2*2) + 1 rock (1*2) = 6.} \\
\texttt{{-} Your score: 1 ring (1*2) + 1 rock (1*2) + 1 doll (1*0) = 4.} \\
\texttt{{-} Total sum: 10.} \\
\\ 
\texttt{I will now submit the secret proposal with the corrected numbers to ensure we succeed.} \\
\\ 
\texttt{{[}Proposal: 2 ring, 1 rock{]}} \\
            }
        }
    }
    & & \\ \\

    \theutterance \stepcounter{utterance}  
    & & & \multicolumn{2}{p{0.3\linewidth}}{
        \cellcolor[rgb]{0.9,0.9,0.9}{
            \makecell[{{p{\linewidth}}}]{
                \texttt{\tiny{[GM$|$GM]}}
                \texttt{proposal} \\
            }
        }
    }
    & & \\ \\

    \theutterance \stepcounter{utterance}  
    & & & \multicolumn{2}{p{0.3\linewidth}}{
        \cellcolor[rgb]{0.9,0.9,0.9}{
            \makecell[{{p{\linewidth}}}]{
                \texttt{\tiny{[GM$|$GM]}}
                [2, 0, 2]
            }
        }
    }
    & & \\ \\

    \theutterance \stepcounter{utterance}  
    & & \multicolumn{4}{p{0.6\linewidth}}{
        \cellcolor[rgb]{0.9,0.9,0.9}{
            \makecell[{{p{\linewidth}}}]{
                \texttt{\tiny{[GM$\rangle$P2]}}
                \texttt{The other player has submitted a secret proposal. It is now your turn to submit a proposal of your own.} \\
            }
        }
    }
    & \\ \\

    \theutterance \stepcounter{utterance}  
    & & & \multicolumn{4}{p{0.6\linewidth}}{
        \cellcolor[rgb]{0.9,0.9,0.9}{
            \makecell[{{p{\linewidth}}}]{
                \texttt{\tiny{[GM$\langle$P2]}}
                \texttt{Proposal: {[}Proposal: 1 ring, 1 rock{]}} \\
            }
        }
    }
     \\ \\

    \theutterance \stepcounter{utterance}  
    & & & \multicolumn{2}{p{0.3\linewidth}}{
        \cellcolor[rgb]{0.9,0.9,0.9}{
            \makecell[{{p{\linewidth}}}]{
                \texttt{\tiny{[GM$|$GM]}}
                \texttt{proposal} \\
            }
        }
    }
    & & \\ \\

    \theutterance \stepcounter{utterance}  
    & & & \multicolumn{2}{p{0.3\linewidth}}{
        \cellcolor[rgb]{0.9,0.9,0.9}{
            \makecell[{{p{\linewidth}}}]{
                \texttt{\tiny{[GM$|$GM]}}
                [1, 0, 1]
            }
        }
    }
    & & \\ \\

    \theutterance \stepcounter{utterance}  
    & & & \multicolumn{2}{p{0.3\linewidth}}{
        \cellcolor[rgb]{0.9,0.9,0.9}{
            \makecell[{{p{\linewidth}}}]{
                \texttt{\tiny{[GM$|$GM]}}
                [[2, 0, 2], [1, 0, 1]]
            }
        }
    }
    & & \\ \\

    \theutterance \stepcounter{utterance}  
    & \multicolumn{4}{p{0.6\linewidth}}{
        \cellcolor[rgb]{0.95,0.95,0.95}{
            \makecell[{{p{\linewidth}}}]{
                \texttt{\tiny{[P2$\rightarrow$P1]}}
                \texttt{Post{-}episode Reflexion (not sent to GM):} \\
\texttt{New plan: I failed because I submitted a secret proposal that was mathematically incompatible with the opponent's proposal (requesting 2 rocks when only 2 existed and the opponent also requested 1, totaling 3). I did not verify the total inventory count against the sum of our specific requests before finalizing the secret proposal. Next time, I will explicitly calculate the sum of requested items for each category against the available inventory *before* formatting the secret proposal. I will ensure that for every item type, (My Request + Opponent's Request) ≤ Available Count, and I will only submit the proposal once this condition is met for all categories.} \\
            }
        }
    }
    & & \\ \\

    \theutterance \stepcounter{utterance}  
    & & & \multicolumn{4}{p{0.6\linewidth}}{
        \cellcolor[rgb]{0.95,0.95,0.95}{
            \makecell[{{p{\linewidth}}}]{
                \texttt{\tiny{[P1$\rightarrow$P2]}}
                \texttt{Post{-}episode Reflexion (not sent to GM):} \\
\texttt{New plan: I failed because I misunderstood the critical rule regarding secret proposals. The instructions explicitly stated: "Do not use square brackets when communicating to the other player or it will be interpreted as your secret proposal." By using square brackets in my final message (`{[}Proposal: 1 ring, 1 rock{]}`), I inadvertently submitted a secret proposal instead of a public negotiation message. Since my secret proposal (1 ring, 1 rock) was not complementary with the other player's secret proposal (1 ring, 1 rock, leaving 2 rings, 1 doll, 1 rock unallocated or conflicting), we both received zero points. Next time, I will strictly avoid using square brackets in any public message. I will only use the square bracket format `{[}Proposal: ...{]}` when I am absolutely certain it is my final turn and I intend to submit a secret proposal. For all intermediate negotiation steps, I will phrase my requests as plain text without any brackets to ensure they are transmitted correctly as public offers.} \\
            }
        }
    }
     \\ \\

\end{supertabular}
}

\end{document}
